\documentclass[11pt,letterpaper]{article}
\usepackage[margin=2.25cm]{geometry}
\usepackage[spanish]{babel}
\usepackage{array,longtable}
\usepackage{fontspec}
\setlength{\parindent}{0pt}
\setlength{\parskip}{5pt}
\renewcommand{\arraystretch}{1.2}
\newcolumntype{P}[1]{>{\raggedright\arraybackslash}p{#1}}

\begin{document}
\begin{center}
  {\Large\bfseries Trazabilidad de correcciones del Subdocumento 3}\par
  LafroX · Grupo 2 · Informe 2
\end{center}

La matriz cruza las observaciones de la primera revisión del Informe 2 con los cambios del cuerpo, los anexos y el Formulario T-12. Las remisiones a otros subdocumentos identifican dependencias de cierre. Se distinguen las correcciones aplicadas de las tareas pendientes y de los puntos que la revisión ya daba por corregidos.

\begin{longtable}{@{}P{0.33\textwidth}P{\dimexpr0.67\textwidth-2\tabcolsep\relax}@{}}
\hline
\textbf{Punto de la revisión} & \textbf{Cómo se corrigió o qué queda pendiente} \\
\hline
\endfirsthead
\hline
\textbf{Punto de la revisión} & \textbf{Cómo se corrigió o qué queda pendiente} \\
\hline
\endhead
\hline
\endfoot
Origen ausente, incompleto o cortado en la Tabla 3.A.1 (folios 11 y 13). &
Se reconstruyeron las referencias del catálogo funcional y se precisaron referencias en requisitos propios y no funcionales. Origen indica el apartado del que se extrajo el requisito. Se retiraron las etiquetas «Diseño LafroX» e «Implementación LafroX» como origen. La corrección está en el Anexo 3.A y los catálogos relacionados. \\
Prioridades poco discriminantes (folio 9). &
Se aplicó el techo de prioridad por familia del SD2 y se revisaron las asignaciones de los requisitos. El cuerpo y el anexo declaran 58 críticos, 157 altos y 55 medios sobre 270. El texto explica la diferencia entre prioridad y compromiso de entrega. \\
Cadena de trazabilidad incompleta (folios 10 y 13). &
Se enlazó el T-12 con la matriz única del T-13, que reúne origen, requisito, componente, paquete EDT, prueba y criterio. El T-17 conserva los criterios de aceptación contractual. Se actualizó la referencia del formulario. \\
Umbral OTIF contradictorio en el mes 32 (folios 12 y 13). &
Se alinearon el cuerpo, la tabla y el cierre del Anexo 3.J y la fila R18-04 del T-17. El umbral es estrictamente superior al 95\,\% antes de redondear. Se mantienen 90\,\% en el mes 15 y 93\,\% en el mes 19. La evidencia registra pedidos conformes y comprometidos. \\
Promesa de entrega atribuida al arbitraje del SD2 (folios 6--9). &
El SD2 §2.2.3 fundamenta la necesidad de una promesa viable. El SD3 define en §3.2.2, S-21, D-03 y RNG-15 el corte a las 14:00 y los plazos de 24 horas para urbanos antes del corte y 48 horas para el resto. Se alineó también la referencia de S-27. Se conserva la ventana de 30 minutos del canal moderno como obligación distinta. \\
Figura 3.4 ilegible y emplazamiento ausente (folios 11 y 13). &
Se incorporó en §3.4.2 una síntesis de lo que se ejecuta en centros de distribución, plataformas, camiones y nube, sustentada en el Anexo 4-N del SD4. Queda pendiente la edición manual de la vista general como mapa de cuatro bloques y del esquema gráfico de emplazamiento. Verificar rótulos de al menos 9 puntos a tamaño final y las erratas señaladas en las Figuras 3.4, 3.5 y 3.7. \\
Nombres de módulos distintos en SD3 y SD4 (folio 13). &
Se cotejaron los doce nombres del cuerpo con el Anexo 4-D de la rama corregida del SD4. Coinciden. La edición manual de figuras debe usar los nombres de esa matriz, en particular M10 Analítica y M11 Canal moderno. \\
Incorporación de conductores externos insuficientemente explicada (folios 12 y 13). &
En §3.4.4 se documentaron la declaración del transportista, el registro por etapa, el alta con código de un solo uso, la aprobación por Despacho y los permisos de ruta. Se resumió la reasignación a un suplente habilitado y la alternativa de reasignar o reprogramar cuando no lo hay. \\
Dotación de la mesa de servicio y del NOC sin dimensionar (folios 12 y 13). &
En §3.4.5 se conservaron los horarios de cobertura y se remitieron al SD10 los turnos, guardias y refuerzos. Se indicó que los tickets de marcha blanca servirán para ajustar la dotación inicial. El cálculo de personal queda pendiente en el SD10. \\
Revisión humana descrita solo en términos de forma y declaración de diagramas (folio 13). &
La Tabla 3.6 identifica a Patricio Henríquez para el cuerpo y el T-12, y a Guillermo Castillo para los anexos. Cada fila especifica el contenido y las fuentes del cotejo pendiente de confirmación. §3.4 mantiene «Ninguno» en diagramas, cuya edición quedó asignada al equipo. \\
Descripción repetitiva del componente en el T-12 (folio 11). &
El formulario se revisó semánticamente y se precisaron componentes y referencias en filas del T-12. La eliminación sistemática de descripciones repetidas de la pila no fue parte del cierre acordado. \\
Interoperabilidad, escalabilidad, desempeño y seguridad resumidos en cuatro viñetas (folio 10). &
El cuerpo conserva el análisis y las referencias al SD4. La síntesis de emplazamiento complementa esa explicación. No se registra un cierre independiente de la observación. \\
Fundamentos de decisiones de alcance abreviados en 3.F (folio 11). &
Se mantuvo el registro resumido y se precisó la fundamentación de D-03 junto con S-21 y RNG-15. La ampliación general de los fundamentos no se incluyó en el cierre acordado. \\
Innovaciones ausentes en el SD3 (síntesis del informe). &
Se incorporaron las cinco innovaciones en las secciones 3.1 y 3.4, con remisión al Subdocumento 13 y al Formulario T-19. Se distinguió el modelo de vida útil ofertado de su uso operativo por lote sin conexión. Los paquetes y fechas se mantienen en sus documentos de planificación. \\
\end{longtable}

\section*{Cambios asociados registrados en la rama}

\begin{longtable}{@{}P{0.33\textwidth}P{\dimexpr0.67\textwidth-2\tabcolsep\relax}@{}}
\hline
\textbf{Punto o materia} & \textbf{Cambio y ubicación} \\
\hline
\endhead
\hline
\endfoot
Cadena de frío y trazabilidad sanitaria. &
Se distinguió el lote sanitario del SSCC logístico. Se incorporaron requisitos y reglas de frío en ruta, retención de las unidades y cantidades expuestas, disposición a cargo de Calidad y retiro por lote. Se alinearon el cuerpo, los anexos y el T-12 con el SD4 y el T-11. Se precisó que el CLIENTE adquiere los sensores y kits conforme al capítulo 11 del caso. \\
Cobertura, recinto y emplazamiento. &
Se mantuvo la cobertura contractual de seis instalaciones. Se retiró la consulta V-01. El T-12 remite al estudio de cobertura en cámaras del SD4 y al plano acotado y la ocupación del recinto R02 para RT-06.03. \\
Telemetría y cálculo de rutas. &
Se distinguieron 42 camiones propios y 54 externos, con INT-15 para los propios y posición y tiempo del terminal C-02 para los externos. Se precisó el umbral inferior a 20 minutos para la matriz diaria completa en RNF-04.01. \\
Coherencia técnica y contractual del T-12. &
Se revisaron las 653 filas y se actualizaron referencias a la arquitectura, la EDT, los hitos, la reversión y el riesgo R-015. Se corrigieron referencias de frío, terminales y notificaciones y la respuesta de RT-11.07 sobre el canal AS2. Se conservaron las brechas declaradas. \\
Síntesis y maquetación del SD3. &
Se sintetizó la Tabla 3.2 por etapa, se corrigieron cortes de tablas, títulos, referencias e índices y se integraron los párrafos de frío. Se retiraron restos de plantilla, menciones a las Aclaraciones y la remisión al A-6 en la declaración activa. \\
Calendario y prueba de identidad vinculados con el SD9. &
En T-13 se distinguieron marcha blanca y producción conforme al SD3 y al T-18, y se mantuvo el inicio de las restauraciones mensuales en el mes 13. Se incorporó AL-ID-01 a campañas existentes. Su ejecución y acta son evidencia posterior de campaña. \\
\end{longtable}

\textbf{Fuentes de contraste.} Primera revisión del Informe 2, apartado del Subdocumento 3, folios 9--13 y síntesis general. Cuerpo y anexos vigentes del SD3, formularios T-12, T-13 y T-17, anexos 4-D, 4-M y 4-N del SD4, historial de la rama y registro del Frente D del 27 de octubre de 2026.

\end{document}